\section{Feature Pyramid Network}

\subsection{Motivazione}

Una rete convoluzionale profonda, come ResNet-50, produce una gerarchia di
feature map a diverse risoluzioni spaziali. I layer più profondi hanno una
risoluzione spaziale più bassa, ma contengono informazioni semantiche più
ricche, mentre i layer più superficiali mantengono una maggiore risoluzione
spaziale e una localizzazione più precisa, ma hanno informazioni semantiche
meno forti.

Questo rappresenta un problema per l'object detection, perché gli oggetti
possono avere dimensioni molto diverse. Le feature profonde sono utili per
riconoscere \textit{cosa} è presente nell'immagine, ma sono meno precise nel
dire \textit{dove} si trova. Al contrario, le feature superficiali sono più
precise spazialmente, ma hanno una rappresentazione semantica più debole.

La Feature Pyramid Network (FPN) serve quindi a combinare queste due
informazioni. L'idea è prendere le feature semantiche dei layer profondi e
propagarle verso i layer a maggiore risoluzione, combinandole con le feature
superficiali. In questo modo si ottengono feature map a diverse scale che
mantengono sia una buona informazione semantica sia una buona precisione
spaziale.\footnote{Tsung-Yi Lin et al., \textit{Feature Pyramid Networks for
Object Detection}, CVPR 2017.}

\subsection{Struttura della FPN}

La FPN originale utilizza le feature ottenute dagli ultimi blocchi residuali
di ResNet, indicate come $C_2$, $C_3$, $C_4$ e $C_5$. Queste feature hanno
rispettivamente stride $4$, $8$, $16$ e $32$ rispetto all'immagine di input.

La costruzione della piramide parte dalla feature più profonda $C_5$. Questa
viene trasformata tramite una convoluzione $1\times1$ e costituisce il punto
di partenza del percorso \textit{top-down}. La feature viene poi
sovracampionata di un fattore 2 e sommata alla corrispondente feature
proveniente dal backbone, dopo che anch'essa è stata trasformata tramite una
convoluzione $1\times1$.

Il procedimento viene ripetuto progressivamente verso i livelli a maggiore
risoluzione:

\[
C_5 \rightarrow P_5
\]

\[
C_4 + \operatorname{Upsample}(P_5) \rightarrow P_4
\]

\[
C_3 + \operatorname{Upsample}(P_4) \rightarrow P_3
\]

\[
C_2 + \operatorname{Upsample}(P_3) \rightarrow P_2
\]

La convoluzione $1\times1$ viene utilizzata per portare le feature del
backbone allo stesso numero di canali, mentre l'upsampling permette di
allineare la risoluzione spaziale delle feature prima della somma.

Dopo la fusione viene applicata una convoluzione $3\times3$ per ottenere la
feature finale di ciascun livello. Nel paper originale tutti i livelli della
piramide hanno 256 canali.\footnote{La FPN originale utilizza il percorso
\textit{top-down} con \textit{lateral connections} e produce i livelli
$P_2$--$P_5$.}

\subsection{FPN utilizzata nel progetto}

Nel nostro progetto la FPN è leggermente diversa da quella originale, perché
è stata adattata al detector utilizzato successivamente.

La ResNet-50 fornisce le feature:

\[
C_3,\ C_4,\ C_5
\]

con stride rispettivamente:

\[
8,\ 16,\ 32.
\]

A differenza della FPN originale, $C_2$ non viene utilizzata. La costruzione
del percorso top-down rimane però la stessa:

\[
C_5 \rightarrow P_5
\]

\[
C_4 + \operatorname{Upsample}(P_5) \rightarrow P_4
\]

\[
C_3 + \operatorname{Upsample}(P_4) \rightarrow P_3.
\]

Anche in questo caso le feature provenienti dal backbone vengono prima
proiettate tramite convoluzioni $1\times1$ a 256 canali. Dopo la somma viene
applicata una convoluzione $3\times3$.

La differenza principale è che il nostro modello non si ferma a $P_3$, $P_4$
e $P_5$, ma estende ulteriormente la piramide:

\[
P_5 \rightarrow P_6 \rightarrow P_7.
\]

$P_6$ viene ottenuta tramite una convoluzione $3\times3$ con stride 2 applicata
a $P_5$, mentre $P_7$ viene ottenuta nello stesso modo partendo da $P_6$.

La piramide finale è quindi:

\[
P_3,\ P_4,\ P_5,\ P_6,\ P_7
\]

con stride:

\[
8,\ 16,\ 32,\ 64,\ 128.
\]

Questa scelta segue la struttura del detector descritto nel paper sulla
pneumonia, che utilizza cinque livelli della feature pyramid con stride
$8$, $16$, $32$, $64$ e $128$.\footnote{Linghua Wu et al.,
\textit{Pneumonia detection based on RSNA dataset and anchor-free deep
learning detector}, Scientific Reports, 2024.}

In sintesi, la FPN permette di avere feature a più scale, combinando
informazioni semantiche profonde con informazioni spaziali più precise dei
livelli superficiali. Nel nostro caso, questi livelli vengono poi utilizzati
dal detector FCOS per effettuare le predizioni a scale differenti.